\uuid{3GrS}
\exo7id{6184}
\titre{exo7 6184}
\auteur{queffelec}
\organisation{exo7}
\datecreate{2011-10-16}
\isIndication{false}
\isCorrection{false}
\chapitre{Compacité}
\sousChapitre{Compacité}
\cpge{mpsi-mp2i.raisonnement-vocabulaire-ensembliste.c, pcsi.raisonnement-vocabulaire-ensembliste.d, ptsi.raisonnement-vocabulaire-ensembliste.d, mpsi-mp2i.limites-continuite.c, pcsi.limites-continuite.c, ptsi.limites-continuite.c, bcpst1.vocabulaire-applications, bcpst1.limites-continuite.c, mp-mpi.topologie-espaces-normes.e, pc.espaces-vectoriels-normes.e, psi.espaces-vectoriels-normes.e, pt.fonctions-deux-trois-variables.b, mp-mpi.topologie-espaces-normes.h}
\module{Topologie}
\niveau{L3}
\difficulte{}

\contenu{
\question{
Soit $f$ une surjection continue de $\Rr^2$ sur $\Rr$. On va montrer que
l'image réciproque de tout point est non bornée. On raisonne par l'absurde : 

 Sinon, il existe $a\in \Rr$ et un disque fermé $ D$ du plan tel que
$f^{-1}(\{a\})\subset  D$; en étudiant $f(D^c)$ et $f(D)$ montrer que
$f(\Rr^2)$ ne peut être égal à $\Rr$ tout entier.
}
}
